% !TeX root = ../../Book1.tex
\section{\texorpdfstring{\(\mathcal H^d\)}{Hd} 与 Lebesgue 测度}
\label{b1:sec:1503}

% 实际读取：Fremlin2016Measure，作者官方 chap26.tex，264H--I 完整证明及264节末注释。
% 264H的有限次坐标对称化在此完整展开；不调用Brunn--Minkowski或面积公式。
% 264I的球覆盖比较改接本书13.1 Vitali，另由第三章立方体覆盖补全零测遗漏。
% 归一化常数以本书070504已证一维Gauss积分、Tonelli和Gamma定义重新推导。

用直径计算覆盖的代价，是否仍能得到通常的长度、面积和体积？
仅凭平移不变性还不能确定答案中的常数。
关键是一个几何问题：给定直径时，集合的体积至多有多大。
本节先证明球达到最大的体积，再把这一不等式与 Vitali 覆盖结合，
精确比较 \(\mathbb R^d\) 上的 \(d\) 维 Hausdorff 测度与 Lebesgue 测度。

\subsection{等直径问题}
\label{b1:subsec:150301}

沿用第\ref{b1:sec:1303}节的记号
\(\omega_d=m_d\bigl(B(0,1)\bigr)\)，其中 \(d\geq1\)。
球面的 Lebesgue 测度为零，所以使用开单位球或闭单位球得到同一个 \(\omega_d\)。
若一个集合的直径为 \(D\)，任选其中一点为中心，只能直接把它放进半径 \(D\) 的球。
要把体积估计改善到半径 \(D/2\) 的球，不能单靠这种包含关系。

\begin{theorem}[等直径不等式]\label{b1:thm:isodiametric}
设 \(A\subseteq\mathbb R^d\) 有界。则
\begin{equation}\label{b1:eq:isodiametric}
 m_d^*(A)\leq\frac{\omega_d}{2^d}
                  (\operatorname{diam}A)^d.
\end{equation}
空集的直径按零理解。对任意 \(D>0\)，直径为 \(D\) 的闭球达到等号。
\end{theorem}

这就是\term[dengzhijing-budengshi]{等直径不等式}[isodiametric inequality]。
它不要求 \(A\) 可测，因此以后可以用于 Hausdorff 覆盖中任意形状的覆盖集。
不等式比较的是体积，而不是说每个直径为 \(D\) 的集合都能放进半径 \(D/2\) 的球。
例如边长为 \(1\) 的正三角形的三个顶点，直径为 \(1\)，却不能放进半径 \(1/2\) 的圆盘：
若一个这样的圆盘含两个相距 \(1\) 的顶点，它的中心必须是二者中点，
第三个顶点距该中点为 \(\sqrt3/2>1/2\)。

下面按 Fremlin 的《Measure Theory》\cite[264H及第264节注释]{Fremlin2016Measure}，
把集合逐坐标对称化。每一步都保持或增大体积，同时不增加直径。
只需 \(d\) 步，就能得到一个关于原点中心对称的集合。

\subsection{欧氏球的极值性质}
\label{b1:subsec:150302}

对第 \(j\) 个坐标，记反射
\[
 R_j(x_1,\ldots,x_d)
 =(x_1,\ldots,x_{j-1},-x_j,x_{j+1},\ldots,x_d).
\]
以下操作把平行于某条坐标轴的截面移到该轴坐标为零的位置，
并用同长度的中心区间替代它。
这种操作称为\term[steiner-duichenhua]{Steiner 对称化}[Steiner symmetrization]；
这里仅需它沿坐标方向的一个版本。

\begin{lemma}[逐坐标对称化]\label{b1:lem:coordinate-symmetrization}
设 \(d\geq2\)、\(C\subseteq\mathbb R^d\) 有界，并且对每个 \(i\in J\subseteq\{1,\ldots,d\}\) 都有 \(R_i(C)=C\)。
若 \(j\notin J\)，则存在有界 Borel 集 \(D\)，使
\[
 m_d(D)\geq m_d^*(C),\quad
 \operatorname{diam}D\leq\operatorname{diam}C,
 \quad R_i(D)=D\quad(i\in J\cup\{j\}).
\]
\end{lemma}
\begin{proof}
置换坐标不会改变距离或 Lebesgue 测度，故只需证明 \(j=d\) 的情形。
空集取 \(D=\varnothing\)。以下令 \(F=\overline C\)，它是紧集，满足
\[
 \operatorname{diam}F=\operatorname{diam}C,
 \quad m_d(F)\geq m_d^*(C).
\]
反射与取闭包可交换，所以 \(F\) 仍保留所有 \(J\) 中的坐标对称性。
对 \(y\in\mathbb R^{d-1}\)，令
\[
 F_y=\{\,t\in\mathbb R:(y,t)\in F\,\},\quad
 \ell(y)=m_1(F_y),\quad
 D=\{\,(y,t):|t|<\ell(y)/2\,\}.
\]
由截面测度的可测性，\(\ell\) 为 Borel 函数，故 \(D\) 是 Borel 集。
紧集 \(F\) 的投影有界，各个截面的长度也有统一上界，因此 \(D\) 有界。
每个 \(D_y\) 的长度恰为 \(\ell(y)\)，包括 \(\ell(y)=0\) 时的空截面。
Tonelli 定理给出
\[
 m_d(D)=\int_{\mathbb R^{d-1}}\ell(y)\,\mathrm dy=m_d(F).
\]

接着验证直径。任取 \((y,t),(y',t')\in D\)，这两个对应的 \(F\)-截面都非空。
记它们的最小、最大点分别为
\[
 a=\min F_y,\quad b=\max F_y,
 \quad a'=\min F_{y'},\quad b'=\max F_{y'}.
\]
这些端点属于截面，因为截面紧。又有 \(\ell(y)\leq b-a\)、\(\ell(y')\leq b'-a'\)，所以
\[
 \begin{aligned}
 |t-t'|
 &\leq|t|+|t'|
 <\frac{\ell(y)+\ell(y')}{2}\\
 &\leq\frac{(b-a)+(b'-a')}{2}
 =\frac{(b-a')+(b'-a)}2\\
 &\leq\max\{\,|b-a'|,|b'-a|\,\}.
 \end{aligned}
\]
因此可以从 \((b,a')\) 与 \((a,b')\) 中选一对端点 \((u,v)\)，使 \(|u-v|\geq|t-t'|\)。
两点 \((y,u),(y',v)\) 都属于 \(F\)，从而
\[
 |(y,t)-(y',t')|^2
 =|y-y'|^2+|t-t'|^2
 \leq|y-y'|^2+|u-v|^2
 \leq(\operatorname{diam}F)^2.
\]
对 \(D\) 中两点取上确界，就得到所需直径估计。

由 \(D\) 的定义，\(R_d(D)=D\)。
若 \(i\in J\)，则 \(i<d\)，并且 \(F_{R_i y}=F_y\)，这里 \(R_i\) 也表示对应的 \(\mathbb R^{d-1}\) 反射。
故 \(\ell(R_i y)=\ell(y)\)，进而 \((y,t)\in D\) 当且仅当 \((R_i y,t)\in D\)。
这说明以前取得的坐标对称性没有被本次操作破坏。
\end{proof}

\begin{proof}[定理\ref{b1:thm:isodiametric}的证明]
空集及直径为零的集合都满足结论。设 \(A\ne\varnothing\)，并记 \(D_0=\operatorname{diam}A>0\)。
若 \(d=1\)，集合包含于 \([\inf A,\sup A]\)，该区间长度为 \(D_0\)，而 \(\omega_1=2\)，故结论成立。

设 \(d\geq2\)。从 \(A_0=A\) 出发，依次对第 \(1,\ldots,d\) 个坐标应用前一引理，得到 \(A_1,\ldots,A_d\)。
它们满足
\[
 m_d^*(A_0)\leq m_d(A_1)\leq\cdots\leq m_d(A_d),
 \quad\operatorname{diam}A_k\leq D_0,
\]
且 \(A_k\) 关于前 \(k\) 个坐标反射不变。
若某一步所得集合为空，上式已经给出 \(m_d^*(A)=0\)。否则对任意 \(x\in A_d\)，连续反射全部坐标得到 \(-x\in A_d\)，故
\[
 2|x|=|x-(-x)|\leq\operatorname{diam}A_d\leq D_0.
\]
因此 \(A_d\subseteq\overline B(0,D_0/2)\)，由体积的伸缩公式可得
\[
 m_d^*(A)\leq m_d(A_d)
 \leq m_d\bigl(\overline B(0,D_0/2)\bigr)
 =\omega_d(D_0/2)^d.
\]
最后，任意半径 \(D_0/2\) 的闭球都具有直径 \(D_0\) 和上述体积，故确实达到等号。
\end{proof}

证明没有把原集合直接装进同直径的球，而是先改变其形状，再比较体积。
中间取闭包可能增大体积，但不增大直径，也不破坏已有的反射对称性；
这正是能够对任意有界集合使用同一论证的原因。
球达到最佳上界，已经足以确定下面的测度常数，不需要再对全部等号情形作分类。

\subsection{\texorpdfstring{\(\mathcal H^d\)}{Hd} 与 Lebesgue 测度}
\label{b1:subsec:150303}

先使用第\ref{b1:sec:1501}节未归一化的覆盖量 \(h_\delta^d\) 与 \(h^d\)。
这样可以先确定几何给出的常数，再与 \(c_d\) 的 Gamma 积分表达式核对。
所用比较见 Fremlin 的《Measure Theory》\cite[264I]{Fremlin2016Measure}。

有一个细节不能省略：Vitali 定理只覆盖几乎所有点，Hausdorff 定义却要求覆盖全部点。
我们先说明 Lebesgue 零测集能用任意小的直径代价补上。
若 \(m_d^*(N)=0\)，给定 \(\delta,\varepsilon>0\)，由\cref{b1:prop:cube-cover}可取半开立方体 \((Q_i)\) 覆盖 \(N\)，使
\[
 \sum_i m_d(Q_i)<\varepsilon d^{-d/2}.
\]
把边长过大的立方体作有限次等分，总代价不变，可使每个新立方体的直径都小于 \(\delta\)。
边长为 \(a\) 的立方体直径为 \(\sqrt d\,a\)，故细分后的覆盖满足
\[
 \sum_i(\operatorname{diam}Q_i)^d
 =d^{d/2}\sum_i m_d(Q_i)<\varepsilon.
\]
这直接从覆盖定义给出 \(h_\delta^d(N)=0\)，进而 \(h^d(N)=0\)，并未预先使用待证的测度一致性。

\begin{proposition}\label{b1:prop:raw-hausdorff-lebesgue}
对任意 \(E\subseteq\mathbb R^d\) 及每个 \(\delta>0\)，都有
\[
 h_\delta^d(E)=h^d(E)=\frac{2^d}{\omega_d}m_d^*(E).
\]
\end{proposition}
\begin{proof}
先证下界。若 \((E_i)\) 是 \(E\) 的一个直径不超过 \(\delta\) 的覆盖，
则由等直径不等式与外测度的次可加性，
\[
 m_d^*(E)\leq\sum_i m_d^*(E_i)
 \leq\frac{\omega_d}{2^d}\sum_i(\operatorname{diam}E_i)^d.
\]
对覆盖取下确界，得到 \(m_d^*(E)\leq(\omega_d/2^d)h_\delta^d(E)\)。
因此若 \(m_d^*(E)=\infty\)，全部所需等式已经成立。

以下设 \(m_d^*(E)<\infty\)。给定 \(\eta>0\)，由\cref{b1:prop:open-cover-approximation}取开集 \(G\supseteq E\)，使
\(m_d(G)<m_d^*(E)+\eta\)。
考虑全部包含于 \(G\)、直径小于 \(\delta\) 的正半径中心闭球。
它们在 \(E\) 上细覆盖，故\cref{b1:thm:vitali-covering-lebesgue}给出至多可数个两两不交的闭球 \((B_j)\)，使
\[
 N=E\setminus\bigcup_jB_j,
 \quad m_d^*(N)=0.
\]
这些球仍包含于 \(G\)，于是
\[
 \sum_j(\operatorname{diam}B_j)^d
 =\frac{2^d}{\omega_d}\sum_j m_d(B_j)
 \leq\frac{2^d}{\omega_d}m_d(G).
\]
再用本小节开头构造的小立方体覆盖 \(N\)，令其直径代价之和小于 \(\varepsilon\)。
合并两族就真正覆盖了 \(E\)，并且每个覆盖集的直径都小于 \(\delta\)。因此
\[
 h_\delta^d(E)
 \leq\frac{2^d}{\omega_d}\bigl(m_d^*(E)+\eta\bigr)+\varepsilon.
\]
先让 \(\varepsilon,\eta\downarrow0\)，得所需上界。
这对每个 \(\delta>0\) 都成立，再取 \(\delta\downarrow0\) 即得到 \(h^d\) 的公式。
\end{proof}

下界使用任意覆盖集的等直径体积界，上界则主动选择球，因为球恰好达到这个界。
Vitali 选择把球的质量之和控制在开集的体积内；剩余零测集由立方体覆盖补齐。
所以两种方向的估计最终给出了同一个常数。

\subsection{归一化常数}
\label{b1:subsec:150304}

第\ref{b1:sec:1501}节对所有实数 \(s\geq0\) 统一采用
\[
 \mathcal H^s=c_sh^s,\quad
 c_s=\frac{\pi^{s/2}}{2^s\Gamma(1+s/2)}.
\]
在整数维数，这个看似解析的表达式恰好等于单位直径球的体积。

\begin{proposition}[整数维归一化]\label{b1:prop:hausdorff-normalization}
对每个正整数 \(d\)，有
\[
 \omega_d\Gamma(1+d/2)=\pi^{d/2},\quad c_d=\frac{\omega_d}{2^d}.
\]
\end{proposition}
\begin{proof}
第\ref{b1:subsec:070504}小节已用二维换元证明
\(\int_{\mathbb R}e^{-t^2}\,\mathrm dt=\sqrt\pi\)。
由 Tonelli 定理反复分离坐标，得到
\[
 \int_{\mathbb R^d}e^{-|x|^2}\,\mathrm dx
 =\left(\int_{\mathbb R}e^{-t^2}\,\mathrm dt\right)^d
 =\pi^{d/2}.
\]
另一方面，恒等式
\(e^{-|x|^2}=\int_{|x|^2}^\infty e^{-u}\,\mathrm du\)
与 Tonelli 定理给出
\[
 \begin{aligned}
 \int_{\mathbb R^d}e^{-|x|^2}\,\mathrm dx
 &=\int_0^\infty e^{-u}
      m_d\bigl(\{\,x:|x|^2\leq u\,\}\bigr)\,\mathrm du\\
 &=\omega_d\int_0^\infty e^{-u}u^{d/2}\,\mathrm du
 =\omega_d\Gamma(1+d/2).
 \end{aligned}
\]
这里仅使用了球体积的伸缩公式，没有使用高维极坐标或球面的面积公式。
比较两次计算，便得第一式，代入 \(c_d\) 的定义得到第二式。
\end{proof}

\begin{theorem}\label{b1:thm:hausdorff-lebesgue}
采用上述归一化时，对任意 \(E\subseteq\mathbb R^d\)，有
\[
 \mathcal H^d(E)=m_d^*(E).
\]
因此 \(\mathcal H^d\) 与 Lebesgue 测度具有相同的 Carathéodory 可测集，
且在这些集合上完全相等。
\end{theorem}
\begin{proof}
由前一命题及\cref{b1:prop:raw-hausdorff-lebesgue}，
\[
 \mathcal H^d(E)
 =c_dh^d(E)
 =\frac{\omega_d}{2^d}\cdot\frac{2^d}{\omega_d}m_d^*(E)
 =m_d^*(E).
\]
这首先是两个外测度在所有集合上的相等。
将其代入 Carathéodory 可测性判据，可知可测集族也相同，限制后的测度自然一致。
\end{proof}

\begin{corollary}\label{b1:cor:hausdorff-flat-subspace}
设 \(1\leq k\leq d\)，\(T\colon\mathbb R^k\to\mathbb R^d\) 是等距仿射嵌入。
则对任意 \(E\subseteq\mathbb R^k\)，有
\[
 \mathcal H^k_{\mathbb R^d}\bigl(T(E)\bigr)=m_k^*(E).
\]
这里的下标仅用于区分计算 Hausdorff 测度的环境空间。
\end{corollary}
\begin{proof}
映射 \(T\) 及其从像空间到 \(\mathbb R^k\) 的逆映射都保持距离。
由\cref{b1:prop:hausdorff-lipschitz}中的等距不变性与子空间环境独立性，
\[
 \mathcal H^k_{\mathbb R^d}\bigl(T(E)\bigr)
 =\mathcal H^k_{\mathbb R^k}(E).
\]
再在 \(k\) 维欧氏空间应用前一定理，得到右侧等于 \(m_k^*(E)\)。
\end{proof}

因此一条直线上的 \(\mathcal H^1\) 是通常的长度，一个平面上的 \(\mathcal H^2\) 是通常的面积，
不论它们嵌入多高维的欧氏空间。
这些是等距参数化的直接结果；曲面或非等距参数化中的体积变化，需要另外计算。

例如 \(c_1=1\)、\(c_2=\pi/4\)、\(c_3=\pi/6\)，分别把直径代价换算成线段长度、圆盘面积和球体积。
这些数值也可由 Gamma 积分直接核对：分部积分给出
\(\Gamma(t+1)=t\Gamma(t)\)（\(t>0\)），而 \(\Gamma(1)=1\)、
\(\Gamma(1/2)=2\int_0^\infty e^{-v^2}\,\mathrm dv=\sqrt\pi\)。
分部积分先在 \([\varepsilon,R]\) 上进行，再令 \(\varepsilon\downarrow0\)、\(R\uparrow\infty\)，
端点项因 \(t>0\) 与指数衰减而消失。
零维时 \(c_0=1\)，仍得到第\ref{b1:sec:1501}节的计数测度。

归一化不会改变某个集合的 Hausdorff 维数，因为对固定 \(s\)，因子 \(c_s\) 是严格正的有限数。
但涉及长度、面积、体积或密度中的精确常数时，归一化必须保持一致。
比较不同参考书的公式，应同时检查覆盖代价究竟写作 \((\operatorname{diam}E_i)^s\)，
还是已经乘入 \(c_s\)，不能只根据符号 \(\mathcal H^s\) 判断它们数值相同。
